\IfFileExists{stacks-project.cls}{%
\documentclass{stacks-project}
}{%
\documentclass{amsart}
}

% For dealing with references we use the comment environment
\usepackage{verbatim}
\newenvironment{reference}{\comment}{\endcomment}
%\newenvironment{reference}{}{}
\newenvironment{slogan}{\comment}{\endcomment}
\newenvironment{history}{\comment}{\endcomment}

% For commutative diagrams we use Xy-pic
\usepackage[all]{xy}

% We use 2cell for 2-commutative diagrams.
\xyoption{2cell}
\UseAllTwocells

% We use multicol for the list of chapters between chapters
\usepackage{multicol}

% This is generally recommended for better output
\usepackage{lmodern}
\usepackage[T1]{fontenc}

% For cross-file-references
\usepackage{xr-hyper}

% Package for hypertext links:
\usepackage{hyperref}

% For any local file, say "hello.tex" you want to link to please
% use \externaldocument[hello-]{hello}
\externaldocument[introduction-]{introduction}
\externaldocument[conventions-]{conventions}
\externaldocument[sets-]{sets}
\externaldocument[categories-]{categories}
\externaldocument[topology-]{topology}
\externaldocument[sheaves-]{sheaves}
\externaldocument[sites-]{sites}
\externaldocument[stacks-]{stacks}
\externaldocument[fields-]{fields}
\externaldocument[algebra-]{algebra}
\externaldocument[brauer-]{brauer}
\externaldocument[homology-]{homology}
\externaldocument[derived-]{derived}
\externaldocument[simplicial-]{simplicial}
\externaldocument[more-algebra-]{more-algebra}
\externaldocument[smoothing-]{smoothing}
\externaldocument[modules-]{modules}
\externaldocument[sites-modules-]{sites-modules}
\externaldocument[injectives-]{injectives}
\externaldocument[cohomology-]{cohomology}
\externaldocument[sites-cohomology-]{sites-cohomology}
\externaldocument[dga-]{dga}
\externaldocument[dpa-]{dpa}
\externaldocument[sdga-]{sdga}
\externaldocument[hypercovering-]{hypercovering}
\externaldocument[schemes-]{schemes}
\externaldocument[constructions-]{constructions}
\externaldocument[properties-]{properties}
\externaldocument[morphisms-]{morphisms}
\externaldocument[coherent-]{coherent}
\externaldocument[divisors-]{divisors}
\externaldocument[limits-]{limits}
\externaldocument[varieties-]{varieties}
\externaldocument[topologies-]{topologies}
\externaldocument[descent-]{descent}
\externaldocument[perfect-]{perfect}
\externaldocument[more-morphisms-]{more-morphisms}
\externaldocument[flat-]{flat}
\externaldocument[groupoids-]{groupoids}
\externaldocument[more-groupoids-]{more-groupoids}
\externaldocument[etale-]{etale}
\externaldocument[chow-]{chow}
\externaldocument[intersection-]{intersection}
\externaldocument[pic-]{pic}
\externaldocument[weil-]{weil}
\externaldocument[adequate-]{adequate}
\externaldocument[dualizing-]{dualizing}
\externaldocument[duality-]{duality}
\externaldocument[discriminant-]{discriminant}
\externaldocument[derham-]{derham}
\externaldocument[local-cohomology-]{local-cohomology}
\externaldocument[algebraization-]{algebraization}
\externaldocument[curves-]{curves}
\externaldocument[resolve-]{resolve}
\externaldocument[models-]{models}
\externaldocument[functors-]{functors}
\externaldocument[equiv-]{equiv}
\externaldocument[pione-]{pione}
\externaldocument[etale-cohomology-]{etale-cohomology}
\externaldocument[proetale-]{proetale}
\externaldocument[relative-cycles-]{relative-cycles}
\externaldocument[more-etale-]{more-etale}
\externaldocument[trace-]{trace}
\externaldocument[crystalline-]{crystalline}
\externaldocument[spaces-]{spaces}
\externaldocument[spaces-properties-]{spaces-properties}
\externaldocument[spaces-morphisms-]{spaces-morphisms}
\externaldocument[decent-spaces-]{decent-spaces}
\externaldocument[spaces-cohomology-]{spaces-cohomology}
\externaldocument[spaces-limits-]{spaces-limits}
\externaldocument[spaces-divisors-]{spaces-divisors}
\externaldocument[spaces-over-fields-]{spaces-over-fields}
\externaldocument[spaces-topologies-]{spaces-topologies}
\externaldocument[spaces-descent-]{spaces-descent}
\externaldocument[spaces-perfect-]{spaces-perfect}
\externaldocument[spaces-more-morphisms-]{spaces-more-morphisms}
\externaldocument[spaces-flat-]{spaces-flat}
\externaldocument[spaces-groupoids-]{spaces-groupoids}
\externaldocument[spaces-more-groupoids-]{spaces-more-groupoids}
\externaldocument[bootstrap-]{bootstrap}
\externaldocument[spaces-pushouts-]{spaces-pushouts}
\externaldocument[spaces-chow-]{spaces-chow}
\externaldocument[groupoids-quotients-]{groupoids-quotients}
\externaldocument[spaces-more-cohomology-]{spaces-more-cohomology}
\externaldocument[spaces-simplicial-]{spaces-simplicial}
\externaldocument[spaces-duality-]{spaces-duality}
\externaldocument[formal-spaces-]{formal-spaces}
\externaldocument[restricted-]{restricted}
\externaldocument[spaces-resolve-]{spaces-resolve}
\externaldocument[formal-defos-]{formal-defos}
\externaldocument[defos-]{defos}
\externaldocument[cotangent-]{cotangent}
\externaldocument[examples-defos-]{examples-defos}
\externaldocument[algebraic-]{algebraic}
\externaldocument[examples-stacks-]{examples-stacks}
\externaldocument[stacks-sheaves-]{stacks-sheaves}
\externaldocument[criteria-]{criteria}
\externaldocument[artin-]{artin}
\externaldocument[quot-]{quot}
\externaldocument[stacks-properties-]{stacks-properties}
\externaldocument[stacks-morphisms-]{stacks-morphisms}
\externaldocument[stacks-limits-]{stacks-limits}
\externaldocument[stacks-cohomology-]{stacks-cohomology}
\externaldocument[stacks-perfect-]{stacks-perfect}
\externaldocument[stacks-introduction-]{stacks-introduction}
\externaldocument[stacks-more-morphisms-]{stacks-more-morphisms}
\externaldocument[stacks-geometry-]{stacks-geometry}
\externaldocument[moduli-]{moduli}
\externaldocument[moduli-curves-]{moduli-curves}
\externaldocument[examples-]{examples}
\externaldocument[exercises-]{exercises}
\externaldocument[guide-]{guide}
\externaldocument[desirables-]{desirables}
\externaldocument[coding-]{coding}
\externaldocument[obsolete-]{obsolete}
\externaldocument[fdl-]{fdl}
\externaldocument[index-]{index}

% Theorem environments.
%
\theoremstyle{plain}
\newtheorem{theorem}[subsection]{Theorem}
\newtheorem{proposition}[subsection]{Proposition}
\newtheorem{lemma}[subsection]{Lemma}

\theoremstyle{definition}
\newtheorem{definition}[subsection]{Definition}
\newtheorem{example}[subsection]{Example}
\newtheorem{exercise}[subsection]{Exercise}
\newtheorem{situation}[subsection]{Situation}

\theoremstyle{remark}
\newtheorem{remark}[subsection]{Remark}
\newtheorem{remarks}[subsection]{Remarks}

\numberwithin{equation}{subsection}

% Macros
%
\def\lim{\mathop{\mathrm{lim}}\nolimits}
\def\colim{\mathop{\mathrm{colim}}\nolimits}
\def\Spec{\mathop{\mathrm{Spec}}}
\def\Hom{\mathop{\mathrm{Hom}}\nolimits}
\def\Ext{\mathop{\mathrm{Ext}}\nolimits}
\def\SheafHom{\mathop{\mathcal{H}\!\mathit{om}}\nolimits}
\def\SheafExt{\mathop{\mathcal{E}\!\mathit{xt}}\nolimits}
\def\Sch{\mathit{Sch}}
\def\Mor{\mathop{\mathrm{Mor}}\nolimits}
\def\Ob{\mathop{\mathrm{Ob}}\nolimits}
\def\Sh{\mathop{\mathit{Sh}}\nolimits}
\def\NL{\mathop{N\!L}\nolimits}
\def\CH{\mathop{\mathrm{CH}}\nolimits}
\def\proetale{{pro\text{-}\acute{e}tale}}
\def\etale{{\acute{e}tale}}
\def\QCoh{\mathit{QCoh}}
\def\Ker{\mathop{\mathrm{Ker}}}
\def\Im{\mathop{\mathrm{Im}}}
\def\Coker{\mathop{\mathrm{Coker}}}
\def\Coim{\mathop{\mathrm{Coim}}}

% Boxtimes
%
\DeclareMathSymbol{\boxtimes}{\mathbin}{AMSa}{"02}

%
% Macros for moduli stacks/spaces
%
\def\QCohstack{\mathcal{QC}\!\mathit{oh}}
\def\Cohstack{\mathcal{C}\!\mathit{oh}}
\def\Spacesstack{\mathcal{S}\!\mathit{paces}}
\def\Quotfunctor{\mathrm{Quot}}
\def\Hilbfunctor{\mathrm{Hilb}}
\def\Curvesstack{\mathcal{C}\!\mathit{urves}}
\def\Polarizedstack{\mathcal{P}\!\mathit{olarized}}
\def\Complexesstack{\mathcal{C}\!\mathit{omplexes}}
% \Pic is the operator that assigns to X its picard group, usage \Pic(X)
% \Picardstack_{X/B} denotes the Picard stack of X over B
% \Picardfunctor_{X/B} denotes the Picard functor of X over B
\def\Pic{\mathop{\mathrm{Pic}}\nolimits}
\def\Picardstack{\mathcal{P}\!\mathit{ic}}
\def\Picardfunctor{\mathrm{Pic}}
\def\Deformationcategory{\mathcal{D}\!\mathit{ef}}
\usepackage{cleveref}
\crefname{originalreference}{original reference}{original references}
\Crefname{originalreference}{Original reference}{Original references}
\makeatletter
\newlabel{definition-cartesian-derived-modules}{{0D9N}{source}{}{sourcepointer-1}{}}
\newlabel{definition-cartesian-derived-modules@cref}{{[originalreference][0][]0D9N}{[1][0][]source}{}{}{}}
\makeatother
\crefname{originalreference}{original reference}{original references}
\Crefname{originalreference}{Original reference}{Original references}
\makeatletter
\newlabel{lemma-cartesian-objects-derived-modules}{{0D9P}{source}{}{sourcepointer-2}{}}
\newlabel{lemma-cartesian-objects-derived-modules@cref}{{[originalreference][0][]0D9P}{[1][0][]source}{}{}{}}
\makeatother
\crefname{originalreference}{original reference}{original references}
\Crefname{originalreference}{Original reference}{Original references}
\makeatletter
\newlabel{lemma-coh-descent-qcoh-unbounded-for-fppf-hypercovering}{{0DHF}{source}{}{sourcepointer-3}{}}
\newlabel{lemma-coh-descent-qcoh-unbounded-for-fppf-hypercovering@cref}{{[originalreference][0][]0DHF}{[1][0][]source}{}{}{}}
\makeatother
\crefname{originalreference}{original reference}{original references}
\Crefname{originalreference}{Original reference}{Original references}
\makeatletter
\newlabel{lemma-hom-cartesian-modules-derived}{{0D9S}{source}{}{sourcepointer-4}{}}
\newlabel{lemma-hom-cartesian-modules-derived@cref}{{[originalreference][0][]0D9S}{[1][0][]source}{}{}{}}
\makeatother
\crefname{originalreference}{original reference}{original references}
\Crefname{originalreference}{Original reference}{Original references}
\makeatletter
\newlabel{lemma-nullity-cartesian-modules-derived}{{0D9R}{source}{}{sourcepointer-5}{}}
\newlabel{lemma-nullity-cartesian-modules-derived@cref}{{[originalreference][0][]0D9R}{[1][0][]source}{}{}{}}
\makeatother
\crefname{originalsection}{original section}{original sections}
\Crefname{originalsection}{Original section}{Original sections}
\makeatletter
\newlabel{section-simplicial-algebraic-spaces}{{0DE7}{source}{}{sourcepointer-6}{}}
\newlabel{section-simplicial-algebraic-spaces@cref}{{[originalsection][0][]0DE7}{[1][0][]source}{}{}{}}
\makeatother
\title{Negative Ext vanishing and the Hom exact sequence along an fppf hypercovering}
\author{The Stacks Project Authors}
\date{stacks-project:spaces-simplicial.tex:local-snapshot-20261012; GNU FDL 1.2 or later}
\begin{document}
\maketitle
\noindent\textit{Editorial note.} All delivered mathematical text is the author's own, copied from the pinned Stacks Project chapter file sources/stacks-spaces-simplicial/pinned-original-source.tex; no statement or proof step was rewritten, reordered or removed.\par\medskip
\noindent\textit{Editorial note.} The counted proof is the author's complete argument: the cohomological-descent input is card L46153 (lemma-coh-descent-qcoh-unbounded-for-fppf-hypercovering); the nullity lemma used is the modules twin of card L46186 (lemma-nullity-cartesian-objects-derived); the assumed componentwise vanishing includes the i < 0 cases needed by the cited lemmas; no verification is omitted.\par\medskip
\noindent\textit{Editorial note.} The label lemma-fppf-neg-ext-zero-hom names the counted statement, and the proof environment immediately following it in the source is delivered verbatim as its proof.\par\medskip
\noindent\textit{Editorial note.} The counted statement is carried by the author's own lemma environment in the chapter source.\par\medskip
\noindent\textit{Allowed standard prerequisites (editorial).}\begin{itemize}
\item The simplicial systems of the derived category attached to La\textasciicircum{}*K and La\textasciicircum{}*E. Same chapter, Definition \textbackslash{}label\{definition-cartesian-derived-modules\} and \textbackslash{}label\{lemma-cartesian-objects-derived-modules\}. Proof lines 7844-7853: K\_n = La\_n\textasciicircum{}*K and E\_n = La\_n\textasciicircum{}*E are the simplicial systems of the derived category of modules associated to La\textasciicircum{}*K and La\textasciicircum{}*E, with a : U\_etale -> X\_etale as in the simplicial-algebraic-spaces section.
\item Unbounded cohomological descent for fppf hypercoverings on QCoh. Same chapter, \textbackslash{}label\{lemma-coh-descent-qcoh-unbounded-for-fppf-hypercovering\} (card L46153). Proof lines 7855-7860: Hom\_\{O\_X\}(K, E) = Hom\_\{O\_U\}(La\textasciicircum{}*K, La\textasciicircum{}*E).
\item Vanishing and the Hom exact sequence for cartesian modules. Same chapter, \textbackslash{}label\{lemma-nullity-cartesian-modules-derived\} and \textbackslash{}label\{lemma-hom-cartesian-modules-derived\}. Proof lines 7862-7872: the first arrow is injective by the nullity lemma; its image is the kernel of the second by the Hom lemma for cartesian modules.
\end{itemize}\medskip
\ifdefined\hypertarget\hypertarget{sourcepointer-1}{}\fi
\noindent\textit{Source reference.} Original reference 0D9N: sources/stacks-spaces-simplicial/pinned-original-source.tex, source lines 3357--3357. Prerequisite pointer invoked at proof line(s) 7847; it names a standard fact of the same chapter that the counted argument uses, and every step of the counted proof built on it is retained verbatim in the delivered body.\par\medskip
\ifdefined\hypertarget\hypertarget{sourcepointer-2}{}\fi
\noindent\textit{Source reference.} Original reference 0D9P: sources/stacks-spaces-simplicial/pinned-original-source.tex, source lines 3389--3389. Prerequisite pointer invoked at proof line(s) 7849; it names a standard fact of the same chapter that the counted argument uses, and every step of the counted proof built on it is retained verbatim in the delivered body.\par\medskip
\ifdefined\hypertarget\hypertarget{sourcepointer-3}{}\fi
\noindent\textit{Source reference.} Original reference 0DHF: sources/stacks-spaces-simplicial/pinned-original-source.tex, source lines 7676--7676. Prerequisite pointer invoked at proof line(s) 7853; it names a standard fact of the same chapter that the counted argument uses, and every step of the counted proof built on it is retained verbatim in the delivered body.\par\medskip
\ifdefined\hypertarget\hypertarget{sourcepointer-4}{}\fi
\noindent\textit{Source reference.} Original reference 0D9S: sources/stacks-spaces-simplicial/pinned-original-source.tex, source lines 3516--3516. Prerequisite pointer invoked at proof line(s) 7872; it names a standard fact of the same chapter that the counted argument uses, and every step of the counted proof built on it is retained verbatim in the delivered body.\par\medskip
\ifdefined\hypertarget\hypertarget{sourcepointer-5}{}\fi
\noindent\textit{Source reference.} Original reference 0D9R: sources/stacks-spaces-simplicial/pinned-original-source.tex, source lines 3493--3493. Prerequisite pointer invoked at proof line(s) 7870; it names a standard fact of the same chapter that the counted argument uses, and every step of the counted proof built on it is retained verbatim in the delivered body.\par\medskip
\ifdefined\hypertarget\hypertarget{sourcepointer-6}{}\fi
\noindent\textit{Source reference.} Original section 0DE7: sources/stacks-spaces-simplicial/pinned-original-source.tex, source lines 6985--6985. Prerequisite pointer invoked at proof line(s) 7851; it names a standard fact of the same chapter that the counted argument uses, and every step of the counted proof built on it is retained verbatim in the delivered body.\par\medskip
\begin{lemma}
\label{lemma-fppf-neg-ext-zero-hom}
Let $X$ be an algebraic space over a scheme $S$.
Let $K, E \in D_\QCoh(\mathcal{O}_X)$.
Let $a : U \to X$ be an fppf hypercovering.
Assume that for all $n \geq 0$ we have
$$
\Ext_{\mathcal{O}_{U_n}}^i(La_n^*K, La_n^*E) = 0
\text{ for } i < 0
$$
Then we have
\begin{enumerate}
\item $\Ext_{\mathcal{O}_X}^i(K, E) = 0$ for $i < 0$, and
\item there is an exact sequence
$$
0
\to
\Hom_{\mathcal{O}_X}(K, E)
\to
\Hom_{\mathcal{O}_{U_0}}(La_0^*K, La_0^*E)
\to
\Hom_{\mathcal{O}_{U_1}}(La_1^*K, La_1^*E)
$$
\end{enumerate}
\end{lemma}

\begin{proof}
Write $K_n = La_n^*K$ and $E_n = La_n^*E$. Then these are the
simplicial systems of the derived category of modules
(Definition \ref{definition-cartesian-derived-modules})
associated to $La^*K$ and $La^*E$
(Lemma \ref{lemma-cartesian-objects-derived-modules})
where $a : U_\etale \to X_\etale$ is as in
Section \ref{section-simplicial-algebraic-spaces}.
Let us prove (2) first. By
Lemma \ref{lemma-coh-descent-qcoh-unbounded-for-fppf-hypercovering}
we have
$$
\Hom_{\mathcal{O}_X}(K, E) =
\Hom_{\mathcal{O}_U}(La^*K, La^*E)
$$
Thus the sequence looks like this:
$$
0
\to
\Hom_{\mathcal{O}_U}(La^*K, La^*E)
\to
\Hom_{\mathcal{O}_{U_0}}(K_0, E_0)
\to
\Hom_{\mathcal{O}_{U_1}}(K_1, E_1)
$$
The first arrow  is injective by
Lemma \ref{lemma-nullity-cartesian-modules-derived}.
The image of this arrow is the kernel of the second
by Lemma \ref{lemma-hom-cartesian-modules-derived}.
This finishes the proof of (2).
Part (1) follows by applying part (2) with
$K[i]$ and $E$ for $i > 0$.
\end{proof}

\end{document}
